\chapter{Reduced form models of default}
A framework for modeling credit risk. Less intuitive than structural models, but easier to implement. Here, the default probability is modeled by the default intensity process. Useful for modeling of derivatives. Flexible approach.

One has to use models when no replicating portfolio exists, the replicating portfolio consists of illiquid instruments (prices are unreliable), or in the case of multi-name derivatives the transaction costs explode when setting up the replicating portfolio (CDO with more than 100 underlying names).

With reduced form models we model the default of a firm directly. (For structural models we need to observe a lot of information for each firm, all data rarely avaible, very complex (hard to compute derivatives prices), unrealistic short- and long-term spreads). We do not try to tie the fundamentals of the firm to the random default time of a company. Instead the time of default $\tau$ is modeled directly and not in relation to the balance sheet. $\tau$ is simply modeled as the first jump time of some underlying process $N_{t}$, hence ( $\tau$ is a stopping time)

$$
\tau=\inf \left\{t>0 \mid N_{t}>0\right\} .
$$

The approach if the specify $N_{t}$ (or how it jumps) so $\tau$ has the desired properties.\\
We call the information generated by the background process $X_{t}$ the background filtration

$$
\left\{\mathcal{G}_{t}\right\}_{t \geq 0}=\left\{\sigma\left\{X_{s} \mid s \leq t\right\}\right\}_{t \geq 0}
$$

All the non-defaultable processes and the intensity process $\lambda_{t}$ is adapted to $\left\{\mathcal{G}_{t}\right\}_{t \geq 0}$. The full filtration $\mathcal{F}_{t}$ is generated by the background filtration and the filtration generated by $N_{t}$

$$
\left\{\mathcal{F}_{t}\right\}_{t \geq 0}=\left\{\mathcal{G}_{t}\right\}_{t \geq 0} \vee\left\{\sigma\left\{N_{s} \mid s \leq t\right\}\right\}_{t \geq 0}
$$

In the intensity reduced form models, the modeling object is the intensity process $\lambda_{t}$ by which the default process $N_{t}$ jumps. We have looked at three different choices Constant intensity $\lambda$, deterministic intensity function $\lambda_{t}$, and stochastic intensity process driven by $X_{t}, \lambda\left(X_{t}\right)$.

